\documentclass[10pt]{article}
\usepackage[a4paper,margin=0.72in]{geometry}
\usepackage{amsmath,amssymb,microtype}
\usepackage{titlesec}
\usepackage[dvipsnames]{xcolor}
\usepackage[colorlinks=true,linkcolor=MidnightBlue,citecolor=MidnightBlue,urlcolor=MidnightBlue]{hyperref}
\titlespacing*{\section}{0pt}{0.8em}{0.25em}
\newcommand{\Hilb}{\mathbf{Hilb}}
\begin{document}
\begin{center}
{\LARGE\bfseries In $\Hilb$, Frobenius implies the $H^*$ axiom}\par
\vspace{0.25em}
{\large The Hilbert-space branch of arXiv:1011.6123 is answered by arXiv:2003.04149}\par
\vspace{0.65em}
Literature-status identification, lane 12 \quad $\cdot$ \quad 13 August 2026
\end{center}
\vspace{0.4em}

\paragraph{Status.}
\texttt{literature\_already\_answered} for the central category $\Hilb$.
Poinsot proves the required semisimplicity, and in fact proves a stronger
equivalence without the source's specialness axiom.  The source's separate
question for arbitrary monoidal dagger categories is not claimed here.

\section*{The source question}

Abramsky and Heunen consider a bounded comultiplication
$\delta:H\to H\otimes H$ satisfying associativity (A), commutativity (C),
specialness (M), and the Frobenius law (F).  They ask whether these axioms
force their $H^*$ axiom (H), first in $\Hilb$ and then in general monoidal
dagger categories \cite{AH}.  Writing
\[
  \mu=\delta^\dagger:H\otimes H\longrightarrow H,
\]
axiom (M) is $\delta^\dagger\delta=I_H$, equivalently
$\mu\mu^\dagger=I_H$.  Their structure theorem shows that, in $\Hilb$, (H)
is equivalent to semisimplicity; they consequently reformulate the question
as whether a nonzero radical Frobenius algebra exists.

\section*{Decisive later theorem}

Poinsot calls $(H,\mu)$ a commutative Hilbertian Frobenius semigroup when
$\mu$ is a bounded commutative associative multiplication and $\mu^\dagger$
is a morphism of $H$-bimodules.  This is exactly the operator form of (A),
(C), and (F).  Theorem~33 of \cite{P} proves, for every such semigroup, the
equivalence
\[
 \text{semisimple}
 \quad\Longleftrightarrow\quad
 \mu\mu^\dagger\text{ is injective}.
\]
Under (M) the latter operator is the identity.  Hence the algebra is
semisimple, so the equivalence established in arXiv:1011.6123 gives (H).
Equivalently, Corollary~32 of \cite{P} says that a radical commutative
Hilbertian Frobenius semigroup must have $\mu=0$; this cannot satisfy
$\mu\mu^\dagger=I_H$ unless $H=0$.  Thus no nontrivial radical example of
the kind sought by the source exists.

Poinsot explicitly records immediately after Theorem~33 that the special
case answers the main $\Hilb$ question of Abramsky--Heunen.  Proposition~39
goes further: a commutative Hilbertian algebra is a Hilbertian $H^*$-algebra
if and only if it is a Hilbertian Frobenius semigroup, without assuming
specialness.

\section*{Scope}

The argument is operator-theoretic and is stated for Hilbert spaces.  It
settles both equivalent $\Hilb$ formulations in the source---(F)$\Rightarrow$(H)
under (A),(C),(M), and nonexistence of a nontrivial special radical Frobenius
algebra.  It does not by itself establish the corresponding implication in
an arbitrary monoidal dagger category.  A bounded exact-phrase and
axiom-based search found no source justifying such a general promotion.

{\small
\begin{thebibliography}{9}
\bibitem{AH}
S.~Abramsky and C.~Heunen,
\emph{$H^*$-algebras and nonunital Frobenius algebras: first steps in
infinite-dimensional categorical quantum mechanics},
Mathematical Structures in Computer Science 22 (2012), 1--24;
\href{https://arxiv.org/abs/1011.6123}{arXiv:1011.6123}.

\bibitem{P}
L.~Poinsot,
\emph{Hilbertian Frobenius algebras},
Communications in Algebra 50 (2022), 85--112;
\href{https://arxiv.org/abs/2003.04149}{arXiv:2003.04149}.
\end{thebibliography}
}

\end{document}
